\section{Related Work}
\label{sec:related}
% =============================================================

\subsection{Short-term solar irradiance forecasting}

Short-term GHI forecasting methods span a wide spectrum from statistical time-series
models to physically-based NWP and data-driven machine learning
\cite{diagne2013,voyant2017}. For intra-hour and sub-day horizons (0--6~h),
satellite-based and ground-observation methods generally outperform NWP, which
requires several hours of spin-up time \cite{perez2013}.

\textbf{Persistence} (assuming future irradiance equals the most recent observation)
is the standard benchmark for short-term forecasting. It is competitive at very
short horizons (under 30~minutes) when irradiance changes slowly, but degrades
rapidly beyond one hour \cite{lonij2013}.

\textbf{Statistical models} such as ARIMA and its seasonal variant SARIMA have
been applied to GHI series after removing the deterministic clear-sky envelope
\cite{reikard2009,bacher2009}. Working on the clear-sky index $k_t = \text{GHI}/\ghics$
improves stationarity and allows standard time-series machinery at hourly resolution.
Despite their simplicity, seasonal ARIMA models can be competitive at multi-hour
horizons in climates with regular diurnal patterns \cite{pedro2012}.

\subsection{Deep learning with satellite imagery}

CNN-based architectures applied to sequences of geostationary satellite patches
have become the dominant approach for satellite-based GHI forecasting. Early
work showed that cloud-motion vectors extracted from visible channels improve
short-term forecasts \cite{chow2011}. More recent studies feed raw satellite
reflectances directly into CNNs followed by fully-connected regressors or LSTMs
\cite{huertas2018,paletta2021,feng2022}. The encoder-decoder design, where a CNN
extracts per-frame spatial embeddings that are then processed by a recurrent model,
has shown consistent gains over purely spatial or purely temporal baselines
\cite{wang2019}.

Attention-based architectures including Vision Transformers (ViT) and Swin
Transformers have also been applied to satellite-based solar forecasting
\cite{acikgoz2022,liu2023}, achieving competitive performance at the cost of
increased parameter count and data requirements. For moderate dataset sizes such
as those encountered in single-site studies, CNN+LSTM hybrids remain a practical
and well-characterised baseline.

\subsection{Graph neural networks for spatial data}

GraphSAGE \cite{hamilton2017} is an inductive GNN that learns aggregation
functions over local neighbourhoods, enabling generalisation to unseen nodes.
It has been applied to traffic forecasting \cite{li2018}, air quality prediction,
and wind power generation \cite{kazem2021}, where spatial proximity defines the
graph topology. In the context of satellite imagery, representing pixels as graph
nodes allows the model to capture irregular spatial dependencies that periodic
CNN filters may miss, such as cloud boundaries or fractional cloud cover patterns.

Closest to our work, \cite{ji2022} applied a GNN-LSTM model to multi-site solar
irradiance forecasting, treating monitoring stations as nodes. Our approach differs
in that we build the graph within a single satellite patch (pixel level), without
requiring multi-station measurements, which is more practical for single-site
deployment.

\subsection{Forecasting in tropical and Andean climates}

Most GHI forecasting benchmarks are conducted in temperate or semi-arid climates
(Europe, southwestern United States, Australia) \cite{paletta2022,benchmarks2023}.
Tropical Andean sites present distinct challenges: convective cloud formation is
driven by topographic lifting and moisture advection, producing strong intraday
variability concentrated in the afternoon \cite{poveda2004}. Few published studies
address satellite-based forecasting for Colombian or equivalent tropical mountain
conditions, motivating the present work.

\subsection{AI-enabled resilience in cyber-physical energy systems}

Cyber-physical energy systems integrate physical generation, transmission, and
distribution assets with digital control and communication layers, creating a
complex interdependency where operational uncertainty propagates bidirectionally
between the cyber and physical domains \cite{deng2017}. The concept of CPES
resilience encompasses the ability to anticipate, withstand, and recover from
disruptive events, including both deliberate attacks and natural uncertainty
sources such as renewable variability \cite{bhusal2020}.

AI-based methods have been applied to improve CPES resilience through improved
situational awareness, predictive analytics, and adaptive control.
For renewable integration specifically, accurate forecasting reduces the need for
spinning reserves and allows pre-emptive load balancing, directly limiting the
impact of generation uncertainty on grid stability \cite{energyrev2021}.
Uncertainty quantification methods that accompany point forecasts with calibrated
confidence intervals enable operators to make risk-aware dispatch decisions:
when forecasts carry explicit probabilistic guarantees, reserve scheduling can
be optimised to the desired risk level rather than relying on conservative
worst-case margins.

This work contributes to the CPES resilience literature by (i) developing
AI-based spatial deep learning for solar generation situational awareness at
tropical sites underrepresented in existing benchmarks, and (ii) coupling
point forecasters with distribution-free conformal prediction intervals
suitable for integration into risk-aware CPES dispatch frameworks.
